Input to the Multi-Aquifer Well (MAW) Package is read from the file that has type ``MAW6'' in the Name File.  Any number of MAW Packages can be specified for a single groundwater flow model.

\vspace{5mm}
\subsubsection{Structure of Blocks}
\vspace{5mm}

\noindent \textit{FOR EACH SIMULATION}
\lstinputlisting[style=blockdefinition]{./mf6ivar/tex/gwf-maw-options.dat}
\lstinputlisting[style=blockdefinition]{./mf6ivar/tex/gwf-maw-dimensions.dat}
\lstinputlisting[style=blockdefinition]{./mf6ivar/tex/gwf-maw-packagedata.dat}
\lstinputlisting[style=blockdefinition]{./mf6ivar/tex/gwf-maw-connectiondata.dat}
\lstinputlisting[style=blockdefinition]{./mf6ivar/tex/gwf-maw-angledata.dat}
\vspace{5mm}
\noindent \textit{FOR ANY STRESS PERIOD}
\lstinputlisting[style=blockdefinition]{./mf6ivar/tex/gwf-maw-period.dat}
\advancedpackageperioddescription{well}{wells}

\vspace{5mm}
\noindent The CONNECTION\_STATUS setting represents a multi-aquifer well connection that has been sealed off from the well, such as a connection that has been backplugged or cased over. A connection is simulated only if both the multi-aquifer well and the connection are active, so a well with a STATUS of INACTIVE has no flow through any of its connections regardless of the status of the individual connections; the connection settings that were specified previously are used again when the well is made ACTIVE. Connection data must be specified for the maximum number of connections that a well will have at any time during the simulation, and a connection that represents a part of a well that has not yet been drilled should be given a CONNECTION\_STATUS of INACTIVE until it is completed.

A packer that isolates part of a multi-aquifer well from the pumped interval, but leaves the isolated interval open to the aquifer, cannot be represented using CONNECTION\_STATUS because the isolated interval has its own water level. Each multi-aquifer well is simulated using a single head, so a well that is divided by a packer should be represented using a separate multi-aquifer well for each isolated interval.

\vspace{5mm}
\subsubsection{Explanation of Variables}
\begin{description}
\input{./mf6ivar/tex/gwf-maw-desc.tex}
\end{description}

\vspace{5mm}
\subsubsection{Example Input File -- Conductance Calculated using Thiem Equation}
\lstinputlisting[style=inputfile]{./mf6ivar/examples/gwf-maw-example1.dat}
\subsubsection{Example Input File -- Conductance Calculated using Screen Geometry}
\lstinputlisting[style=inputfile]{./mf6ivar/examples/gwf-maw-example2.dat}
\subsubsection{Example Input File -- Flowing Well with Conductance Specified}
\lstinputlisting[style=inputfile]{./mf6ivar/examples/gwf-maw-example3.dat}
\subsubsection{Example Input File -- Non-Vertical (Slanted and Horizontal) Connections}
\lstinputlisting[style=inputfile]{./mf6ivar/examples/gwf-maw-example4.dat}

\vspace{5mm}
\subsubsection{Available observation types}
Multi-Aquifer Well Package observations include well head and all of the terms that contribute to the continuity equation for each multi-aquifer well. Additional LAK Package observations include the conductance for a well-aquifer connection conductance (\texttt{conductance}) and the calculated flowing well-aquifer connection conductance (\texttt{fw-conductance}). The data required for each MAW Package observation type is defined in table~\ref{table:gwf-mawobstype}. Negative and positive values for \texttt{maw} observations represent a loss from or gain to the GWF model, respectively. For all other flow terms, negative and positive values represent a loss from or gain to the MAW package, respectively.

\begin{longtable}{p{2cm} p{2.75cm} p{2cm} p{1.25cm} p{7cm}}
\caption{Available MAW Package observation types} \tabularnewline

\hline
\hline
\textbf{Stress Package} & \textbf{Observation type} & \textbf{ID} & \textbf{ID2} & \textbf{Description} \\
\hline
\endfirsthead

\captionsetup{textformat=simple}
\caption*{\textbf{Table \arabic{table}.}{\quad}Available MAW Package observation types.---Continued} \\

\hline
\hline
\textbf{Stress Package} & \textbf{Observation type} & \textbf{ID} & \textbf{ID2} & \textbf{Description} \\
\hline
\endhead

\hline
\endfoot

\input{../Common/gwf-mawobs.tex}
\label{table:gwf-mawobstype}
\end{longtable}

\vspace{5mm}
\subsubsection{Example Observation Input File}
\lstinputlisting[style=inputfile]{./mf6ivar/examples/gwf-maw-example-obs.dat}
